% Plan: development/outline.md, section 6, item B; round-2 adjudication G5-02
% (moves M1 and M2). Statements and proofs moved verbatim from the r3 sources:
% part (c) of prop:split-affine (sections/04-split.tex) with part (c) of its
% proof and the paragraph after the proof (sections/G-proofs-split.tex), and
% the last statement of prop:split-cellwise with the "Best constants" part of
% its proof and the remark on the cap (G-proofs-split.tex,
% app:splitproofs-cellwise).
\subsection{Two explanatory results on splits}\label{app:split-explanatory}

No certificate uses these results; they explain why the certificates choose their slopes as they do.
We use the path model, the stage residuals $F^\varphi_t$ and the split bounds $B(\varphi;K')$ of \cref{sec:split-framework}, and the conventions of \cref{app:splitproofs-setting}.

\begin{proposition}[forced slopes]\label{prop:split-forced}
Let $\varphi$ be an affine split as in \cref{prop:split-affine}, with slopes $\lambda_t$ and $\lambda_0=\lambda_N=0$.
Suppose $B(\varphi)=\vstar$, $\vstar$ is attained at $z^*$, and each $z^*_t$ is an interior point of $K_t$ at which $F_t$ is differentiable.
Then $\nabla F_t(z^*_t)+(\pi^R_t)^\top\lambda_t-(\pi^L_t)^\top\lambda_{t-1}=0$ for every $t$.
These equations determine $\lambda_1,\dots,\lambda_{N-1}$ one after the other from $z^*$.
\end{proposition}

\begin{proof}
By \cref{prop:split-affine}(b), $z^*_t$ minimizes $F^\varphi_t$ over $K_t$.
It is an interior point at which $F^\varphi_t$ is differentiable, so the gradient of $F^\varphi_t$ vanishes there; this is the stated equation.
Because $\pi^R_t$ selects distinct coordinates, the vector $(\pi^R_t)^\top\lambda_t$ carries $\lambda_{t,c}$ at the coordinate $i_c$ of $z_t$ that separator coordinate $c$ selects, and $0$ elsewhere.
The component $i_c$ of the equation for bag $t$ therefore gives $\lambda_{t,c}=-\partial F_t/\partial z_{t,i_c}(z^*_t)+\bigl[(\pi^L_t)^\top\lambda_{t-1}\bigr]_{i_c}$, which determines $\lambda_t$ from $\lambda_{t-1}$, starting from $\lambda_0=0$.
\end{proof}

Unrolled, the recursion says that the slope of a separator coordinate is minus the sum of the partial derivatives of the costs, at $z^*$, with respect to the copies of the same model variable in the bags up to that separator.
The Karush--Kuhn--Tucker conditions of the copy formulation at $z^*$, with all bags interior, are the same equations with the copy-row multipliers in place of $\lambda$, so the slopes of an exact affine split are those multipliers.
When $K_t$ contains equality rows such as discretized dynamics, its interior is empty and \cref{prop:split-forced} does not apply.
After these rows are eliminated, the same computation at a minimizer whose reduced stage problems are interior gives the discrete costate equations, as in the Lagrange form of Krotov's conditions \citep{krotov1967-sufficient-conditions-for-the-optimality}.

\begin{proposition}[best constants for cellwise slopes]\label{prop:split-cellwise-best}
In the setting of \cref{prop:split-cellwise}, suppose in addition that every $\mathcal P_t$ partitions $\R^{k_t}$, every $b_t$ equals the infimum on its right, and $\mathrm{SP}$ is finite.
Then $\mathrm{SP}$ is the largest value of $B(\varphi;K')$ over the cellwise-affine splits $\varphi_t(s)=\lambda_{t,D}^\top s+c_{t,D}$ ($s\in D\in\mathcal P_t$) with these slopes.
\end{proposition}

\begin{proof}
Below, the pair problem of $(D,D')$ at bag $t$ is the infimum on the right of the inequality that defines $b_t(D,D')$ in \cref{prop:split-cellwise}.
Because the cells partition the separator spaces, the sets $\{z\in K'_t:\pi^L_tz\in D,\ \pi^R_tz\in D'\}$ partition $K'_t$.
On each of them the stage residual of the cellwise-affine split equals the objective of the pair problem of $(D,D')$ plus $c_{t,D'}-c_{t-1,D}$, where $c_{0,\ast}=c_{N,\ast}=0$, and $b_t(D,D')$ is the value of that pair problem.
So the stage infimum is
\[
  m_t(c)=\min_{D\in\mathcal P_{t-1},\,D'\in\mathcal P_t}
  \bigl[b_t(D,D')+c_{t,D'}-c_{t-1,D}\bigr],
  \qquad B(\varphi;K')=\sum_{t=1}^Nm_t(c).
\]
Let $(D^*_t)$ be a shortest path.
Bounding each $m_t(c)$ by its value on the arc $(D^*_{t-1},D^*_t)$, the constants telescope and $B(\varphi;K')\le\mathrm{SP}$ for every choice of constants.

For equality we construct finite potentials.
Let $\dist_t(D)\in\R\cup\{+\infty\}$ be the shortest distance from $\ast$ in layer $0$ to $D$ in layer $t$, so $\dist_0(\ast)=0$, $\dist_N(\ast)=\mathrm{SP}$, and $\dist_t(D')\le \dist_{t-1}(D)+b_t(D,D')$ for all arcs.
Every $\dist_t(D)$ exceeds $-\infty$ because the graph is finite and no arc has length $-\infty$.
Prefixes of a shortest path are shortest, so $\dist_t(D^*_t)=\dist_{t-1}(D^*_{t-1})+b_t(D^*_{t-1},D^*_t)$.
Since $\mathrm{SP}<+\infty$, each layer has a finite arc; let $\beta_t$ be the smallest finite arc length in layer $t$.
Choose a real $C_0\ge0$ so large that $M_t=C_0+\sum_{s\le t}\beta_s$ satisfies $M_t\ge \dist_t(D^*_t)$ for $1\le t<N$ and $M_N\ge\mathrm{SP}$, and put
\[
  p_t(D)=\min\{\dist_t(D),M_t\}\in\R,\qquad c_{t,D}=-p_t(D)\qquad(1\le t<N).
\]
Then $p_0(\ast)=0$ because $C_0\ge0$.
\begin{itemize}
\item For $1\le t<N$ and every finite arc, $p_t(D')\le p_{t-1}(D)+b_t(D,D')$:
  if $p_{t-1}(D)=\dist_{t-1}(D)$, then $p_t(D')\le \dist_t(D')\le \dist_{t-1}(D)+b_t(D,D')$;
  otherwise $p_{t-1}(D)=M_{t-1}$ and $p_t(D')\le M_t=M_{t-1}+\beta_t\le M_{t-1}+b_t(D,D')$.
  So every reduced arc length $b_t(D,D')-p_t(D')+p_{t-1}(D)$ is nonnegative.
  On the arc $(D^*_{t-1},D^*_t)$ it is $0$, since $p_t(D^*_t)=\dist_t(D^*_t)$ for all $t<N$.
  Hence $m_t(c)=0$.
\item In the last layer, $m_N(c)=\min_D[b_N(D,\ast)+p_{N-1}(D)]$.
  Each term is at least $\mathrm{SP}$: either $b_N(D,\ast)+\dist_{N-1}(D)\ge \dist_N(\ast)$, or $b_N(D,\ast)+M_{N-1}\ge M_N\ge\mathrm{SP}$ for finite $b_N(D,\ast)$.
  The term for $D^*_{N-1}$ equals $\mathrm{SP}$.
\end{itemize}
So $B(\varphi;K')=0+\dots+0+\mathrm{SP}=\mathrm{SP}$.
\end{proof}

The cap $M_t$ is needed because a cell that no path from the left reaches has $\dist_t=+\infty$, and the choice $c_{t,D}=-\dist_t(D)$ would then not give a real-valued split.
The assumption that no pair bound equals $-\infty$ keeps the shortest-path recursion well defined.
